\section{Physical interface and the exact calibration bridge}
\label{sec:tpc27-interface-bridge}

We retain the actual opened-row provenance before any Mellin or Fourier
separation.  This is essential only at physical additive frequency
zero; the nonzero spectrum will be separated later.

\subsection{Rows, scales, and coefficients}

Fix \(h\in\Z\setminus\{0\}\), and put
\begin{equation}
 H=\rad|h|,
 \qquad H_0=|h|.
 \label{eq:tpc27-H-H0}
\end{equation}
Let
\begin{equation}
 R=X^{r_0+o(1)},
 \qquad r_0=\frac12-\delta>0,
 \qquad R\le S\le X^{C_S},
 \label{eq:tpc27-RS}
\end{equation}
where all exponent inequalities below have fixed positive margins.  In
one opened slice, rows are indexed by
\begin{equation}
 \alpha=(\ell,d_\alpha),
 \qquad m_\alpha=\ell d_\alpha\asymp Q=LD,
 \label{eq:tpc27-row-index}
\end{equation}
with
\begin{equation}
 \ell\in[L/2,2L]\text{ prime},
 \quad d_\alpha\in[D,2D]\cap[1,R^{1/2}],
 \quad d_\alpha\text{ squarefree},
 \quad(d_\alpha,H)=1.
 \label{eq:tpc27-row-support}
\end{equation}
We assume
\begin{equation}
 QJ\asymp X,
 \qquad L>\max\{2R,2H_0\},
 \qquad L\ge X^{\lambda_0}
 \label{eq:tpc27-row-scales}
\end{equation}
for a fixed \(\lambda_0>0\).  In particular, \(\ell>R\).  On the
primitive orbit \((j,H)=1\), target primitivity gives
\begin{equation}
 (u,m_\alpha H)=1
 \qquad\text{whenever }u\mid m_\alpha j+h
 \text{ and }(j,H)=1.
 \label{eq:tpc27-target-primitivity}
\end{equation}
Since \(H\) is fixed and \(d_\alpha\le R^{1/2}\), we also have
\(d_\alpha H\le R\) for all sufficiently large \(X\).  This is the
finite-model completeness condition used below.

The two physical row coefficients have the form
\begin{equation}
 \gamma_\alpha^{(i)}
 =\mu(d_\alpha)(\log\ell)\omega_D^{(i)}(d_\alpha)
  \psi_L^{(i)}(\ell/L)\zeta_\alpha^{(i)},
 \qquad |\zeta_\alpha^{(i)}|\ll1,
 \qquad i=1,2,
 \label{eq:tpc27-row-provenance}
\end{equation}
where the smooth and fixed-residue data range over fixed bounded
families.  Chebyshev's estimate gives
\begin{align}
 \sum_\alpha|\gamma_\alpha^{(i)}|&\ll Q,
 \label{eq:tpc27-row-l1}\\
 \|\gamma^{(i)}\|_\infty&\ll\log(2L),
 \qquad
 \|\gamma^{(i)}\|_2^2\ll Q\log(2L).
 \label{eq:tpc27-row-l2}
\end{align}
These are the provenance bounds of TPC-25
\citep{WangTPC25}.  The generic row-pair mask satisfies
\begin{equation}
 |\mathfrak m(\alpha,\gamma)|\le1,
 \qquad
 \mathfrak m(\alpha,\gamma)=0
 \quad\text{if }m_\alpha=m_\gamma,
 \label{eq:tpc27-mask}
\end{equation}
and is independent of all divisor and conductor variables.  The second
condition will be used in the large-Ramanujan-modulus tail.

Let \(\Xi_{\alpha,\gamma}(j)\) be a uniformly bounded real periodic
factor of fixed period \(H_0\).  The row--orbit weight
\(\cW_{\alpha,\gamma}(j/J)\) belongs to the factorable physical family
of TPC-25.  It is a uniformly bounded \(C_c^\infty\) family supported
in one fixed compact interval, so its sampled orbit has length
\(O(J)\).  In particular,
\begin{equation}
 \sup_{\alpha,\gamma}
 \left|\int_\R\cW_{\alpha,\gamma}(t)\,dt\right|\ll1,
 \label{eq:tpc27-zero-functional}
\end{equation}
and continuous Mellin separation has finite weighted variation in
every fixed orbit seminorm.  We denote all fixed shift, smoothness,
residue, source, and exponent-margin data by \(\mathscr D\).

\subsection{Finite model, prime density, and row drift}

Put
\begin{equation}
 a(u)=-\mu(u)\log u
 \label{eq:tpc27-a}
\end{equation}
and
\begin{equation}
 \lambda'_R(u)
 =u\sum_{b\le R/u}
 \frac{\mu(ub)\mu(b)}{\varphi(ub)}
 \quad(u\le R),
 \qquad \lambda'_R(u)=0\quad(u>R).
 \label{eq:tpc27-lambda-prime}
\end{equation}
Define
\begin{equation}
 b_R(u)=a(u)-\lambda'_R(u).
 \label{eq:tpc27-bR}
\end{equation}
All three coefficients vanish off squarefree integers.  The finite-model
and prime conditional densities are
\begin{align}
 \rho_{M,H,R}(m_\alpha)
 &=\frac H{\varphi(H)}
   \frac{d_\alpha}{\varphi(d_\alpha)},
 \label{eq:tpc27-model-density}\\
 \rho_{P,H}(m_\alpha)
 &=\frac{m_\alpha H}{\varphi(m_\alpha H)}.
 \label{eq:tpc27-prime-density}
\end{align}
Since \(\ell>R\), the large row prime occurs only in the second
density.  Their difference is
\begin{equation}
 \delta_H(m_\alpha)
 :=\rho_{P,H}(m_\alpha)-\rho_{M,H,R}(m_\alpha)
 =\frac H{\varphi(H)}
  \frac{d_\alpha}{\varphi(d_\alpha)(\ell-1)}
 \ll_h\frac{\log\log(3X)}{L}.
 \label{eq:tpc27-row-drift}
\end{equation}

For \(M,Y\ge1\), write
\begin{align}
 A_M(Y)&=
 \sum_{\substack{u\le Y\\(u,M)=1}}
 \frac{-\mu(u)\log u}{u},
 \label{eq:tpc27-AM}\\
 C_M(Y)&=
 \sum_{\substack{u\le Y\\(u,M)=1}}
 \frac{\mu(u)}u.
 \label{eq:tpc27-CM}
\end{align}
The classical zero-free region for \(\zeta(s)\), in the uniform form
proved and used in TPC-20 and TPC-24, gives for every fixed \(A,C>0\)
\begin{align}
 A_M(Y)&=\frac M{\varphi(M)}
 +O_{A,C}\!\left(
  \frac M{\varphi(M)}(\log Y)^{-A}\right),
 \label{eq:tpc27-uniform-A}\\
 C_M(Y)&\ll_{A,C}
 \frac M{\varphi(M)}(\log Y)^{-A},
 \label{eq:tpc27-uniform-C}
\end{align}
uniformly for \(Y\ge3\) and \(M\le Y^C\)
\citep{IwaniecKowalski2004,WangTPC20,WangTPC24}.  These are proved
theorems, not conjectural distribution hypotheses.

\subsection{Centered and calibrated prefixes}

For \(R\le S\), define
\begin{align}
 Z_{m,S}(j)
 &:=\sum_{\substack{u\le S\\(u,mH)=1}}
 b_R(u)\left(\one_{u\mid mj+h}-\frac1u\right),
 \label{eq:tpc27-Zms}\\
 E_S(m)&:=A_{mH}(S)-\rho_{P,H}(m),
 \label{eq:tpc27-ES}\\
 P_{m,S}(j)&:=
 \sum_{\substack{u\le S\\u\mid mj+h}}b_R(u)-\delta_H(m).
 \label{eq:tpc27-Pms}
\end{align}

\begin{proposition}[Exact centered-to-calibrated bridge]
\label{prop:tpc27-PZE}
For every retained row, every \(R\le S\), and every primitive orbit
point \((j,H)=1\),
\begin{equation}
 \boxed{P_{m,S}(j)=Z_{m,S}(j)+E_S(m).}
 \label{eq:tpc27-PZE}
\end{equation}
Moreover, for every fixed \(A>0\),
\begin{equation}
 \boxed{E_S(m)\ll_{A,\mathscr D}(\log X)^{-A}}
 \label{eq:tpc27-E-bound}
\end{equation}
uniformly in the retained polynomial row and cutoff ranges.
\end{proposition}

\begin{proof}
The finite model has the exact density identity
\begin{equation}
 \sum_{\substack{u\le R\\(u,mH)=1}}
 \frac{\lambda'_R(u)}u=\rho_{M,H,R}(m).
 \label{eq:tpc27-model-density-identity}
\end{equation}
Since \(\lambda'_R\) is supported on \(u\le R\), subtraction from
\(A_{mH}(S)\) gives
\begin{equation}
 \sum_{\substack{u\le S\\(u,mH)=1}}
 \frac{b_R(u)}u
 =\delta_H(m)+E_S(m).
 \label{eq:tpc27-prefix-density}
\end{equation}
Expand the centered indicator in \eqref{eq:tpc27-Zms}, insert
\eqref{eq:tpc27-prefix-density}, and use target primitivity
\eqref{eq:tpc27-target-primitivity}.  This proves
\eqref{eq:tpc27-PZE}.  Estimate \eqref{eq:tpc27-E-bound} is
\eqref{eq:tpc27-uniform-A} with \(M=mH\), since all relevant quantities
are fixed powers of \(X\).  The factor
\(mH/\varphi(mH)\ll_h\log\log(3X)\) is absorbed by first invoking
\eqref{eq:tpc27-uniform-A} with one additional logarithmic power.
\end{proof}

The physical base packet is
\begin{align}
 \cQ_S^{\rm phys}
 :={}&\sum_{\alpha,\gamma}
 \gamma_\alpha^{(1)}\gamma_\gamma^{(2)}
 \mathfrak m(\alpha,\gamma)
 \notag\\
 &\quad\times
 \sum_{\substack{j\in\Z\\(j,H)=1}}
 \Xi_{\alpha,\gamma}(j)
 \cW_{\alpha,\gamma}(j/J)
 P_{m_\alpha,S}(j)P_{m_\gamma,S}(j).
 \label{eq:tpc27-physical-base}
\end{align}
The natural scale is
\begin{equation}
 JQ^2\asymp XQ.
 \label{eq:tpc27-natural-scale}
\end{equation}
By \cref{prop:tpc27-PZE}, the integrand in
\eqref{eq:tpc27-physical-base} has the exact four-term decomposition
\begin{equation}
 P_{m,S}P_{n,S}
 =Z_{m,S}Z_{n,S}
 +E_S(m)Z_{n,S}
 +Z_{m,S}E_S(n)
 +E_S(m)E_S(n).
 \label{eq:tpc27-ZZ-EZ-ZE-EE}
\end{equation}
The first term contains the new additive-zero problem.  The last three
will be closed directly in \cref{sec:tpc27-nonzero-scalar}.
